%% ============================================================================
%%  formulario_t11_datacenters.tex
%%
%%  Conversión a LaTeX, sin cambios de contenido, de
%%  04_arquitectura_fisica/4.3/formulario-t11-datacenters.md.
%%  Documento independiente (un formulario por archivo): se compila por sí
%%  solo con  latexmk formulario_t11_datacenters.tex  y no se monta en
%%  main.tex ni en contenido.tex.
%% ============================================================================
\documentclass[carta,firma]{lafrox}

\makeatletter
\@removefromreset{table}{chapter}
\makeatother
\renewcommand{\thetable}{\arabic{table}}

\datoslafrox{
  documento  = Propuesta Técnica,
  subtitulo  = Formulario T-11 --- Data Center Primaria y Data Center Secundario,
  alcance    = Subdocumento 4 --- Arquitectura de la solución,
  formulario = {Formulario T-11},
  fecha      = 05 de octubre de 2026,
  version    = 2.0,
}

\begin{document}

\chapter*{Formulario T-11 --- Fila Data Center Primaria y Data Center Secundario (borrador alineado a v5-431 y v6-432)}

\nota{\textbf{Alcance:} solo las filas de los ítems d) y e) del Subdocumento 4.2. El resto del formulario (cómputo, red, dispositivos de terreno, nube, software) ya vive en \texttt{partes/15\_anexo\_t11.tex} y no se repite aquí.
Columnas oficiales: Componente | Producto / servicio ofertado | Ubicación / Lugar | Cantidad | Justificación (Formulario T-11 de las Bases Administrativas). Sin precios (Art. 50.2).
La tabla de sala técnica del anexo actual \textbf{debe reemplazarse} por esta (valores alineados al cálculo RT-06.11 de la sección 4.3.1).}

\section*{Data Center Primaria --- infraestructura de la sala técnica secundaria de Talca}

\begin{tablalafrox}{Data Center Primaria --- infraestructura de la sala técnica secundaria de Talca}{tab:t11dc-primaria}%
  {F{0.15} Y F{0.17} G{0.10} Y}%
  {\cab{Componente} & \cab{Producto / servicio ofertado} & \cab{Ubicación / Lugar} & \cab{Cantidad} & \cab{Justificación}}
  UPS & Modular de 10 kVA, doble conversión on-line, configuración N+1 y bypass de mantenimiento & Zona de UPS y baterías (plano RT-06.03) & 1 & Dimensionado sobre la carga TI de diseño de 7,0 kW y $\approx$ 7,4 kVA aparentes (FP 0,95), al 80 \% de utilización (resulta $\approx$ 9,2 kVA $\rightarrow$ 10 kVA; RT-06.11; Tabla 4.3-1); autonomía $\geq$ 30 min a plena carga (RT-06.07) \\
  Generador & Grupo electrógeno de 15 kVA con estanque para 24 h continuas y contrato de reabastecimiento declarado & Exterior, junto a la acometida (zona de generadores del plano RT-06.03) & 1 & Sostiene la carga total del sitio $\approx$ 10,4 kW (TI + climatización + apoyo) durante $\geq$ 24 h (RT-06.08) \\
  Transferencia automática & Transferencia red--generador con tablero de protecciones & Acometida, accesible desde el exterior & 1 & Encadena empalme $\rightarrow$ tablero $\rightarrow$ transferencia $\rightarrow$ UPS $\rightarrow$ PDU $\rightarrow$ equipo sin interrupción (RT-06.07); se prueba cada mes con carga real y admite termografía (RT-06.10) \\
  Climatización & 2 unidades de precisión de 10 kW ($\approx$ 34.000 BTU/h) en configuración N+1 & Zona de climatización, contigua a la sala de servidores y comunicaciones & 2 & Cubren la carga térmica sensible de diseño $\approx$ 7,6 kW incluso con la pérdida de la unidad mayor, dentro de los rangos del fabricante (RT-06.13; RT-06.11; Tabla 4.3-2) \\
  Distribución eléctrica de racks & 2 PDU verticales A/B por gabinete, con medidor & Racks R01 y R02 de la sala & 4 & Dos caminos eléctricos por gabinete (RT-08.04); el medidor es el denominador del PUE (RT-06.10; RT-15.04) \\
  Detección de incendio & Detección temprana por aspiración de aire con tecnología láser, tipo AnaLASER o equivalente & Sala de servidores y comunicaciones & 1 & Detecta antes de que exista llama visible (RT-06.16); integrado al monitoreo en línea (RT-06.19) \\
  Extinción & Agente limpio tipo FM-200 o equivalente, con aprobación UL, botón de aborto y extintores portátiles habilitados & Sala de servidores y comunicaciones & 1 & Extinción automática sobre equipos energizados conforme a RT-06.17 y extintores con mantención vigente (RT-06.18). No se cita edición específica de norma externa \\
  Control de acceso & Biometría facial con AFIS como respaldo, esclusa antipassback con nueva verificación, bitácora electrónica y estación de enrolamiento interna y externa & Acceso del recinto técnico & 1 & RT-06.20 a RT-06.23; bitácora auditable conservada $\geq$ 5 años (RT-06.21; RT-16.10) \\
  Videovigilancia & Cámaras IP con imágenes en línea $\geq$ 30 días y respaldo recuperable & Recinto y perímetro & Según plano RT-06.03 & Cubre el acceso y el perímetro del recinto, integrada al monitoreo (RT-06.24) \\
  Gabinetes & R01 servidores (3 nodos 2U en U18--U23, NAS 2U en U1--U2, KVM 1U en U24, paneles OM4/Cat6A en U40--U42) y R02 comunicaciones (bandeja del operador con ONT de fibra y router LTE, 2U en U31--U32; 2 firewalls 1U en U33--U34; switch de gestión 1U en U35; organizadores 1U en U36 y U39; 2 conmutadores de núcleo 1U en U37--U38; ODF y paneles en U40--U42), 42U & Sala de servidores y comunicaciones & 2 & Racks de servidores independientes de los racks de comunicaciones (RT-06.05); ocupación proyectada R01 y R02 de 12U/42U (29 \%) cada uno, con 30U libres por rack para el margen de crecimiento a tres años, coherente con la reserva eléctrica del 20 \% de la Tabla 4.3-1 (Figura 4.3-2) \\
  Piso técnico y cableado & Cableado Cat6A y fibra OM4 certificados por enlace, sobre piso técnico & Sala & 1 & Cumple la norma por enlace (RT-06.04) con recorridos verificados para 10 GbE \\
\end{tablalafrox}

\section*{Data Center Secundario --- gabinete de borde del CD Concepción}

\begin{tablalafrox}{Data Center Secundario --- gabinete de borde del CD Concepción}{tab:t11dc-secundario}%
  {F{0.15} Y F{0.17} G{0.10} Y}%
  {\cab{Componente} & \cab{Producto / servicio ofertado} & \cab{Ubicación / Lugar} & \cab{Cantidad} & \cab{Justificación}}
  Gabinete de borde & Gabinete industrializado con alimentación protegida (UPS + respaldo) para la autonomía declarada & CD Concepción & 1 & Sitio de recuperación del dominio on-premise, a $\approx$ 250--400 km de Talca (RT-07.02); opera de forma autónoma todos los días (RT-06.01 del caso; numeral 6.1 de las Transversales) \\
  Climatización del borde & Climatización de precisión acorde al equipamiento del borde & CD Concepción, gabinete de borde & 1 & Dimensionada al equipamiento real del borde conforme a la tipología del numeral 6.1 \\
  Control de acceso y monitoreo & Acceso controlado y monitoreo remoto integrado al NOC del sitio & CD Concepción, gabinete de borde & 1 & Aplica la tipología de borde operacional del numeral 6.1 con supervisión remota del sitio \\
  Servidor de borde & Dell R250 o HPE DL20 Gen11, 32 GB, RAID 10 sobre las 4 bahías (sin repuesto en caliente) & CD Concepción, gabinete de borde & 1 & Sostiene el WMS con 24 h de operación autónoma; promoción controlada ante contingencia de Talca (RNF-02.01; RT-07.02). Fuente única declarada como punto único de falla aceptado (RT-02.11) \\
  Perímetro del borde --- firewall & Firewall de borde de grado empresarial & CD Concepción, gabinete de borde & 1 & Punto único de falla aceptado (RT-02.11); es el Customer Gateway del túnel hacia AWS y conmuta entre fibra y LTE en menos de 30 s (RT-03.17) \\
  Perímetro del borde --- switch & Switch de borde de grado empresarial & CD Concepción, gabinete de borde & 1 & Punto único de falla aceptado (RT-02.11); sostiene la red de bodega del sitio secundario durante 24 h de autonomía, sin depender del principal \\
\end{tablalafrox}

\nota{Los componentes de la \textbf{réplica en nube} (\texttt{us-east-1}) ya están declarados en el anexo actual: N-04 (réplica reducida ECS), N-05 (Aurora Global Database), N-06 (Global Tables), N-11 (AWS Backup + S3 Object Lock + replicación de S3) y D-01 (VPC/Route 53) --- no requieren cambios.}

\section*{Qué corregir en \texttt{partes/15\_anexo\_t11.tex} (tabla de sala técnica)}

\begin{itemize}
  \item \textbf{D1 UPS:} 6 kVA $\rightarrow$ \textbf{10 kVA modular N+1} (cálculo RT-06.11).
  \item \textbf{D2 Generador:} 12 kVA $\rightarrow$ \textbf{15 kVA} (sostiene $\approx$ 10,4 kW del sitio, RT-06.08).
  \item \textbf{D4 Climatización:} 12.000 BTU/h $\rightarrow$ \textbf{10 kW ($\approx$ 34.000 BTU/h) por unidad, N+1} (carga térmica $\approx$ 7,6 kW, RT-06.13). Retirar la cita ASHRAE (norma externa).
  \item \textbf{D7 Extinción:} retirar la edición de NFPA (``NFPA 75 y 2001'') $\rightarrow$ solo RT-06.17.
  \item \textbf{D11 Gabinetes:} 4 racks (R01--R04) $\rightarrow$ \textbf{2 racks: R01 servidores, R02 comunicaciones}, con margen de crecimiento del 20 \% dentro de los racks (RT-06.05, Figura 4.3-2).
  \item \textbf{D5 PDU:} 4 gabinetes / 8 PDU $\rightarrow$ \textbf{2 racks / 4 PDU}.
  \item \textbf{C7 Switch de gestión:} la consola KVM se declara en ``gabinete R02'' $\rightarrow$ \textbf{R01} (la Figura 4.3-2 la ubica en R01, U24). El switch de gestión permanece en R02.
  \item \textbf{C8 Distribución horizontal:} ``gabinete R03'' $\rightarrow$ \textbf{racks R01 y R02} (los paneles de fibra OM4 y cobre Cat6A están en ambos gabinetes según la Figura 4.3-2; R03 no existe en el diseño vigente).
  \item Verificar D8/D9/D10 coherentes con el plano RT-06.03 (Figura 4.3-1).
\end{itemize}

\end{document}
